\usepackage[french]{babel}
\usepackage{wallpaper}
\usepackage{changepage}
\usepackage[explicit,clearempty]{titlesec}
\usepackage{fancyhdr}
\usepackage{fontspec}
\usepackage{xcolor}
\usepackage{tikz}
\usetikzlibrary{shapes, positioning, shadows}
\usepackage{enumitem}
\usepackage{multicol}
\usepackage{etoolbox}
\usepackage[final]{pdfpages}
\usepackage{xpatch}
\usepackage[a4paper, top=2.5cm, bottom=14mm, left=2.5cm, right=2cm]{geometry}
\usepackage{makeidx}
\makeindex

% Configuration des polices 
\babelfont{rm}{Baskerville}
\babelfont{sf}{Arial}
\newfontfamily\titreFont{Amilya}[Path=./Amilya/, Extension=.ttf, UprightFont=*]

% Définition des couleurs
\definecolor{MyStrongBlue}{RGB}{0, 68, 148}
\definecolor[named]{MyStrongBlue}{HTML}{304fac}
\definecolor[named]{GreenTea}{HTML}{CAE8A2}
\definecolor[named]{MilkTea}{HTML}{C5A16F}
\definecolor[named]{TitleBox}{HTML}{C5A16F}
\definecolor[named]{AuthorName}{HTML}{f2dcbd}
\definecolor[named]{ChapterBox}{HTML}{C5A16F}
\definecolor[named]{ChapterTextBox}{HTML}{c78a36}

% Color-coded recipe elements
\definecolor{ingredientcolor}{RGB}{219, 112, 147}  % Pink
\definecolor{cookwarecolor}{RGB}{100, 149, 237}     % Blue
\definecolor{timercolor}{RGB}{255, 140, 0}          % Orange

% Custom commands for recipe elements
\newcommand{\ingredient}[1]{\textcolor{ingredientcolor}{\textbf{#1}}\index{#1}}
\newcommand{\cookware}[1]{\textcolor{cookwarecolor}{\textbf{#1}}}
\newcommand{\timer}[1]{\textcolor{timercolor}{\textbf{#1}}}

%% Command to hold chapter illustration image
\newcommand\chapterillustration{}
 
%% Define how the chapter title is printed
\titleformat{\chapter}{}{}{0pt}{
  %% Background image at top of page
  \ThisULCornerWallPaper{1}{\chapterillustration}
  %% Draw a semi-transparent rectangle across the top
  \tikz[overlay,remember picture]
    \fill[GreenTea,opacity=0]
    (current page.south west) rectangle 
    ([yshift=-105mm] current page.north east);
  %% Check if on an odd or even page
  \strictpagecheck\checkoddpage
  \ifoddpage{
    \begin{tikzpicture}[overlay,remember picture]
    \node[fill=GreenTea!80!black,text=white,
      font=\fontspec{Marker Felt}\Huge\bfseries, 
      inner ysep=12pt, inner xsep=20pt,
      rounded rectangle,anchor=east, 
      xshift=-40mm,yshift=-10mm] 
      at (current page.north east) {#1};
    \end{tikzpicture}
  }
  \else {
    \begin{tikzpicture}[overlay,remember picture]
    \node[fill=ChapterTextBox!80!black,text=white,
      font=\fontspec{Marker Felt}\Huge\bfseries,
      inner sep=12pt, inner xsep=20pt,
      rounded rectangle,anchor=west,
      xshift=20mm,yshift=-30mm] 
      at ( current page.north west) {#1};
    \end{tikzpicture}
  }
  \fi
  \clearpage
  \thispagestyle{empty}
}

%% Define how the chapter title is printed
\titleformat{name=\chapter,numberless}{}{}{0pt}{
  %% Draw a semi-transparent rectangle across the top
  \tikz[overlay,remember picture]
    \fill[ChapterBox,opacity=.7]
    (current page.north west) rectangle 
    ([yshift=-3cm] current page.north east);
  %% Check if on an odd or even page
  \strictpagecheck\checkoddpage
  \ifoddpage{
    \begin{tikzpicture}[overlay,remember picture]
    \node[fill=ChapterTextBox!80!black,text=white,
      font=\fontspec{Marker Felt}\Huge\bfseries,
      inner ysep=12pt, inner xsep=20pt,
      rounded rectangle,anchor=east, 
      xshift=-20mm,yshift=-30mm] 
      at (current page.north east) {#1};
    \end{tikzpicture}
  }
  \else {
    \begin{tikzpicture}[overlay,remember picture]
    \node[fill=ChapterTextBox!80!black,text=white,
      font=\Huge\bfseries,
      inner sep=12pt, inner xsep=20pt,
      rounded rectangle,anchor=west,
      xshift=20mm,yshift=-30mm] 
      at ( current page.north west) {#1};
    \end{tikzpicture}
  }
  \fi
}
\titlespacing*{\chapter}{0pt}{0pt}{135mm}
\titlespacing*{name=\chapter,numberless}{0pt}{0pt}{50mm}

%% Define how the part title is printed
\titleformat{\part}[display]
  {\thispagestyle{empty}}
  {}
  {0pt}
  {
    \begin{tikzpicture}[remember picture, overlay]
        \node[
            fill=TitleBox, 
            text=Cornsilk, 
            rounded corners=16pt,
            drop shadow={
                shadow xshift=5pt, 
                shadow yshift=-5pt, 
                opacity=0.4,
                black
            },
            inner sep=36pt,
            align=center,
            font=\bfseries
        ] at (current page.center) {
            {\titreFont\fontsize{80pt}{60pt}\selectfont #1}
        };
    \end{tikzpicture}
  }
\titlespacing*{\part}{0pt}{0pt}{0pt}

% Espacement : {commande}{gauche}{avant}{après}[dessous]
% Le 3ème argument (avant) contrôle la hauteur du début du texte. 
% On utilise une valeur faible (ex: 0pt ou même négative) pour remonter le titre.
\titlespacing*{\chapter}{0pt}{-30pt}{20pt}

%% Set the uniform width of the colour box displaying the page number in footer
\newlength\pagenumwidth
\settowidth{\pagenumwidth}{99}

%% Define style of page number colour box
\tikzset{pagefooter/.style={
    anchor=base,font=\sffamily\bfseries\small,
    text=white,fill=MyStrongBlue!80!black,text centered,
    text depth=2mm,text width=2mm+\pagenumwidth
}}

%%%%%%%%%%
%%% Re-define running headers on non-chapter pages
%%%%%%%%%%
\fancypagestyle{headings}{%
  \fancyhf{}   % Clear all headers and footers first
  %% Right headers on odd pages
  \fancyhead[RO]{%
    \begin{tikzpicture}[remember picture,overlay]
    \fill[MilkTea!25!white] (current page.north east) rectangle (current page.south west);
    \fill[white, rounded corners] ([xshift=-10mm,yshift=-20mm]current page.north east) rectangle ([xshift=15mm,yshift=17mm]current page.south west);
    \end{tikzpicture}
    \begin{tikzpicture}[xshift=-.75\baselineskip,yshift=.25\baselineskip,remember picture, overlay,fill=GreenTea,draw=GreenTea]\fill circle(3pt);\draw[semithick](0,0) -- (current page.west |- 0,0);\end{tikzpicture} \sffamily\itshape\small\nouppercase{\rightmark}
  }

  %% Left headers on even pages
  \fancyhead[LE]{%
    \begin{tikzpicture}[remember picture,overlay]
    \fill[MilkTea!25!white] (current page.north east) rectangle (current page.south west);
    \fill[white, rounded corners] ([xshift=-15mm,yshift=-20mm]current page.north east) rectangle ([xshift=10mm,yshift=17mm]current page.south west);
    \end{tikzpicture}
    \sffamily\itshape\small\nouppercase{\leftmark}\ 
    \begin{tikzpicture}[xshift=.5\baselineskip,yshift=.25\baselineskip,remember picture, overlay,fill=GreenTea,draw=GreenTea]\fill (0,0) circle (3pt); \draw[semithick](0,0) -- (current page.east |- 0,0 );\end{tikzpicture}
  }

  \fancyfoot[RO,LE]{\tikz[baseline]\node[pagefooter]{\thepage};}
  \renewcommand{\headrulewidth}{0pt}
  \renewcommand{\footrulewidth}{0pt}
}

%%%%%%%%%%
%%% Re-define running headers on chapter pages
%%%%%%%%%%
\fancypagestyle{plain}{%
  \fancyhf{}
  \fancyfoot[RO,LE]{\tikz[baseline]\node[pagefooter]{\thepage};}
  \renewcommand{\headrulewidth}{0pt}
  \renewcommand{\footrulewidth}{0pt}
}

% tirets dans la table des matières (index)
\makeatletter
\renewcommand*\l@chapter[2]{%
  \ifnum \c@tocdepth >\m@ne
    \addpenalty{-\@highpenalty}%
    \vskip 1.0em \@plus\p@
    \setlength\@tempdima{1.5em}%
    \begingroup
      \parindent \z@ \rightskip \@pnumwidth
      \parfillskip -\@pnumwidth
      \leavevmode \bfseries
      \advance\leftskip\@tempdima
      \hskip -20pt
      #1\nobreak\ 
       \leaders\hbox{$\m@th
        \mkern \@dotsep mu\hbox{.}\mkern \@dotsep
        mu$}\hfil\nobreak\hb@xt@\@pnumwidth{\hss #2}\par
      \penalty\@highpenalty
    \endgroup
  \fi}
\makeatother
